\section{Results}
\label{sec:results}

\subsection{Closed-Set Discrimination and the Generalization Gap}
\label{sec:closedsetresults}

Before evaluating compositional generalization, an auxiliary closed-set
experiment was conducted to determine whether the dataset was intrinsically
difficult to discriminate when all diagnostic classes were represented
during training. The 24-class closed-set classifier achieved an accuracy of
0.9132, a balanced accuracy of 0.9134, and a macro-$F_1$ score of 0.9119,
as summarized in Table~\ref{tab:closedset}.

\begin{table}[!t]
\caption{Closed-set multiclass diagnostic performance with 95\% bootstrap
confidence intervals.}
\label{tab:closedset}
\centering
\begin{tabular}{l c c}
\toprule
Metric & Estimate & 95\% CI \\
\midrule
Accuracy          & 0.9132 & $[0.8854, 0.9375]$ \\
Balanced accuracy & 0.9134 & $[0.8856, 0.9377]$ \\
Macro-$F_1$       & 0.9119 & $[0.8826, 0.9366]$ \\
\bottomrule
\end{tabular}
\end{table}

These results demonstrate strong conventional closed-set discriminability
when all diagnostic classes are represented during training. Consequently,
the substantially lower constituent-level performance observed under the
crossed protocols is unlikely to arise solely from poor baseline
separability; rather, it emerges under the additional requirements of
composition, compound-context, and operating-condition generalization.

\subsection{Crossed Composition--Context--Condition Generalization}
\label{sec:crossedresults}

The proposed crossed evaluation protocol was applied to all 108 unique
physical compound-fault runs. Each target compound was evaluated under
recombination, one-sided compound-context zero-shot, and two-sided
compound-context zero-shot regimes, with each regime further separated into
seen- and unseen-operating-condition settings. In the one-sided regime,
both directional removals were evaluated because either constituent could
be the mechanism deprived of compound-context exposure.
Table~\ref{tab:crossed} summarizes the principal constituent-level results.

\begin{table*}[!t]
\caption{Constituent-level generalization performance under crossed
composition-context and operating-condition regimes. EM denotes exact
constituent-set recovery, $R_{\mathrm{const}}$ denotes constituent recall,
FA/run denotes the mean number of incorrectly added inactive constituents
per compound prediction row, and $R_{\mathrm{any}}$ and $R_{\mathrm{all}}$
denote recovery of at least one and all true constituents, respectively.}
\label{tab:crossed}
\centering
\begin{tabular}{l l c c c c c c c}
\toprule
Condition & Context regime & Active Macro-$F_1$ & Micro-$F_1$ & EM &
$R_{\mathrm{const}}$ & FA/run & $R_{\mathrm{any}}$ & $R_{\mathrm{all}}$ \\
\midrule
Seen   & One-sided     & 0.5890 & 0.6581 & 0.1204 & 0.6505 & 0.6528 & 0.9259 & 0.3750 \\
Seen   & Recombination & 0.6177 & 0.6775 & 0.1389 & 0.6759 & 0.6389 & 0.9352 & 0.4167 \\
Seen   & Two-sided     & 0.5994 & 0.6620 & 0.1111 & 0.6620 & 0.6759 & 0.9444 & 0.3796 \\
Unseen & One-sided     & 0.5761 & 0.6452 & 0.1111 & 0.6481 & 0.7222 & 0.9398 & 0.3565 \\
Unseen & Recombination & 0.6177 & 0.6865 & 0.1574 & 0.6944 & 0.6574 & 0.9444 & 0.4444 \\
Unseen & Two-sided     & 0.5588 & 0.6311 & 0.0833 & 0.6296 & 0.7315 & 0.9352 & 0.3241 \\
\bottomrule
\end{tabular}
\end{table*}

\begin{table*}[!t]
\caption{Input-independent constant-predictor references on the 108
compound runs. These references quantify performance obtainable from
constituent prevalence alone.}
\label{tab:trivial}
\centering
\begin{tabular}{l c c c c c c}
\toprule
Constant prediction & Micro-$F_1$ & EM & $R_{\mathrm{const}}$ & FA/run &
$R_{\mathrm{any}}$ & $R_{\mathrm{all}}$ \\
\midrule
Inner $+$ outer                & 0.5556 & 0.1111 & 0.5556 & 0.8889 & 1.0000 & 0.1111 \\
Inner $+$ outer $+$ broken bar & 0.5333 & 0.0000 & 0.6667 & 1.6667 & 1.0000 & 0.3333 \\
Inner only                     & 0.3704 & 0.0000 & 0.2778 & 0.4444 & 0.5556 & 0.0000 \\
\bottomrule
\end{tabular}
\end{table*}

When both constituents retained prior exposure to other compound contexts,
previously unseen pairings remained partially recoverable. Under the
recombination protocol, micro-$F_1$ was 0.6775 under seen operating
conditions and 0.6865 when the operating condition was also unseen.
Constituent recall showed a similar pattern, reaching 0.6759 and 0.6944,
respectively.

Exact constituent-set recovery remained limited, ranging from 0.0833 to
0.1574 across the six crossed regimes. However, the trivial references in
Table~\ref{tab:trivial} show that some apparently favorable
partial-recovery metrics are strongly influenced by constituent prevalence.
In particular, the constant inner-plus-outer predictor achieves
$R_{\mathrm{any}} = 1.0$ without using the input signals. Consequently,
$R_{\mathrm{any}}$ is retained only as a descriptive measure and is not
interpreted as evidence of diagnostic skill.

Constituent recall alone is not sufficient to establish diagnostic skill in
this setting. The constant inner-plus-outer-plus-broken-bar predictor
reaches $R_{\mathrm{const}} = 0.6667$ without using the input signals, but
does so at the cost of 1.6667 false additions per run and zero exact
matches. Accordingly, constituent recall is interpreted jointly with false
additions, micro-$F_1$, and exact-set recovery throughout the remainder of
the analysis.

The locked logistic-regression detector achieved micro-$F_1$ values of
0.631--0.686 across the six crossed regimes, exceeding the best
constant-predictor micro-$F_1$ of 0.556. At the same time, FA/run remained
between 0.639 and 0.732, substantially below the 1.667 false additions
produced by the three-constituent constant reference and also below the
0.889 value of the constant inner-plus-outer predictor. Under
recombination, $R_{\mathrm{all}}$ reached 0.4167 and 0.4444 for seen and
unseen conditions, respectively, compared with 0.1111 for the
two-constituent constant reference. These results indicate nontrivial
constituent discrimination beyond prevalence-driven constant prediction,
while the low exact-match rates show that reliable exact decomposition of
unseen compound faults remains unresolved.

\subsection{Paired Effects of Compound-Context Deprivation}
\label{sec:pairedeffects}

Paired run-cluster bootstrap analysis was used to determine which
differences between the context regimes were supported beyond variation
across physical runs. Figure~\ref{fig:contrasts} summarizes the
prespecified contrasts for constituent recall and micro-$F_1$.

\begin{figure*}[!t]
\centering
\includegraphics[width=\textwidth]{fig2_contrasts}
\caption{Paired run-cluster bootstrap differences between
composition-context regimes. Points denote paired differences in
constituent recall and micro-$F_1$, and horizontal bars denote 95\%
bootstrap confidence intervals. Positive values favor the first regime in
each contrast.}
\label{fig:contrasts}
\end{figure*}

Under seen operating conditions, the differences between recombination and
the context-zero-shot regimes were limited. For constituent recall, the
recombination--one-sided contrast was 0.0255 with a 95\% CI of
$[0.0000, 0.0509]$, whereas the recombination--two-sided contrast was
0.0139 with a 95\% CI of $[-0.0186, 0.0509]$. The corresponding
micro-$F_1$ contrasts also included zero.

A different pattern was observed when the target operating condition was
withheld. Relative to one-sided context zero-shot, recombination improved
constituent recall by 0.0463 (95\% CI: $[0.0231, 0.0694]$) and micro-$F_1$
by 0.0413 (95\% CI: $[0.0223, 0.0636]$). Relative to two-sided context
zero-shot, the recombination advantage was 0.0648 for constituent recall
(95\% CI: $[0.0278, 0.1065]$) and 0.0554 for micro-$F_1$ (95\% CI:
$[0.0232, 0.0895]$).

The one-sided--two-sided differences remained unresolved for both metrics
under unseen operating conditions. A formal difference-of-differences
analysis further showed that the penalty associated with complete
compound-context deprivation increased under operating-condition shift. For
recombination versus two-sided context zero-shot, the interaction
difference was 0.0400 for micro-$F_1$ (95\% CI: $[0.0088, 0.0731]$), 0.0509
for constituent recall (95\% CI: $[0.0185, 0.0880]$), and 0.0833 for
all-constituents recovery (95\% CI: $[0.0278, 0.1481]$). In contrast, the
recombination-versus-one-sided interaction was not resolved for
micro-$F_1$ or constituent recall, although the interaction for
all-constituents recovery was positive at 0.0463 (95\% CI:
$[0.0046, 0.1019]$). These findings indicate that operating-condition shift
amplifies the penalty of complete compound-context deprivation rather than
producing a uniform interaction across all context regimes.

Directional decomposition of the one-sided regime revealed substantial
mechanism dependence. When broken-bar compound context was withheld,
constituent recall remained high at 0.979 under seen conditions and 0.896
under unseen conditions. In contrast, withholding winding compound context
reduced recall to 0.438 and 0.479, respectively, while withholding
dynamic-eccentricity context yielded a recall of 0.500 under both condition
settings. Each of the 108 physical compound runs contributed exactly two
directional one-sided predictions, and all bootstrap uncertainty estimates
remained clustered by physical run.

\subsection{Sensitivity to Decision-Threshold Objective}
\label{sec:thresholdsensitivity}

The limited exact-match performance and non-negligible false-addition
counts motivated a post-hoc sensitivity analysis of the inner-CV threshold
objective. The original $F_1$-based thresholds were retained as the locked
primary analysis. Alternative $F_{0.5}$ and balanced-accuracy objectives
were evaluated only to characterize the precision--recall operating-point
trade-off and were not used for post-hoc model selection.

Precision-weighted $F_{0.5}$ calibration substantially reduced false
additions. Across the six crossed regimes, FA/run decreased from
0.639--0.731 under the locked $F_1$ thresholds to approximately
0.222--0.389, while exact-match recovery increased from 0.083--0.157 to
0.157--0.287. This gain was accompanied by lower constituent recall, which
decreased to approximately 0.583--0.644. Conversely, balanced-accuracy
thresholding increased constituent recall but also increased false
additions.

These results show that exact compound decomposition is strongly affected
by the operating point of the constituent detectors. Precision-oriented
thresholding suppresses over-diagnosis and improves exact-set recovery,
whereas recall-oriented thresholding recovers more true constituents at the
cost of additional constituent activations. As shown by the
classifier-family robustness analysis below, absolute decomposition quality
is also model-dependent. Consequently, the low exact-match performance of
the locked linear probe should not be interpreted as an intrinsic upper
bound imposed by the crossed protocol itself.

\subsection{Sensitivity to Single-Fault Severity Merging}
\label{sec:severityresults}

Removing low-severity single-fault support did not produce a uniform change
in constituent recall, exact-match recovery, or all-constituent recovery
across the crossed regimes. Most paired differences for these metrics had
95\% confidence intervals spanning zero, indicating that the principal
generalization findings were not driven by pooling low- and high-severity
single-fault examples.

The clearest differences were observed in false additions. Under seen
conditions, removing low-severity support increased FA/run from 0.653 to
0.755 in the one-sided regime ($\Delta = +0.102$, 95\% CI:
$[0.046, 0.162]$) and from 0.639 to 0.722 under recombination
($\Delta = +0.083$, 95\% CI: $[0.019, 0.157]$). Under unseen conditions,
the two-sided regime similarly increased from 0.731 to 0.843 FA/run
($\Delta = +0.111$, 95\% CI: $[0.037, 0.185]$). For seen-condition
recombination, micro-$F_1$ decreased from 0.677 to 0.651 when low-severity
support was removed ($\Delta = -0.026$, 95\% CI: $[-0.052, -0.002]$).
Other micro-$F_1$ differences were unresolved. Overall, the sensitivity
analysis indicates that severity merging does not account for the principal
crossed-generalization pattern and, in several regimes, the additional
low-severity support reduced over-diagnosis rather than artificially
inflating constituent recovery.

\subsection{Robustness Across Classifier Families}
\label{sec:familyresults}

The complete crossed protocol was additionally evaluated with
random-forest, RBF-SVM, and neural classifiers while retaining the same
constituent representation, curated feature families, outer evaluation
splits, and inner-CV threshold-calibration principles.
Table~\ref{tab:families} summarizes descriptive performance averaged across
the six crossed scenarios. The locked logistic-regression results remain
the primary analysis; the alternative classifiers are reported only as
robustness references.

\begin{table}[!t]
\caption{Descriptive classifier-family robustness averaged across the six
crossed evaluation scenarios. LR denotes logistic regression and MLP
denotes the three-hidden-layer feed-forward neural baseline.}
\label{tab:families}
\centering
\begin{tabular}{l c c c c}
\toprule
Model & Micro-$F_1$ & EM & $R_{\mathrm{const}}$ & FA/run \\
\midrule
LR      & 0.660 & 0.120 & 0.660 & 0.680 \\
RF      & 0.744 & 0.282 & 0.669 & 0.256 \\
RBF-SVM & 0.692 & 0.157 & 0.664 & 0.509 \\
MLP     & 0.737 & 0.310 & 0.659 & 0.256 \\
\bottomrule
\end{tabular}
\end{table}

Absolute decomposition performance was classifier dependent. Random forest
achieved the highest mean micro-$F_1$ (0.744), whereas the neural baseline
achieved the highest mean exact-match rate (0.310). Both random forest and
the neural baseline reduced mean FA/run to 0.256, compared with 0.680 for
the locked logistic-regression probe. In contrast, mean constituent recall
varied only from 0.659 to 0.669 across the four classifiers. Thus,
classifier family primarily affected the precision of constituent-set
formation and the number of additional activations rather than the overall
fraction of true constituents recovered.

The qualitative recombination advantage was nevertheless retained across
all four classifier families. Under unseen operating conditions, the
recombination--two-sided difference in micro-$F_1$ was 0.055 for logistic
regression, 0.051 for random forest, 0.032 for RBF-SVM, and 0.023 for the
neural baseline. The corresponding constituent-recall differences were
0.065, 0.065, 0.069, and 0.032, respectively. These descriptive results
indicate that the adverse effect of complete compound-context deprivation
is not unique to the locked linear classifier.

\subsection{Oracle Detectability and Cross-Partner Constituent Transfer}
\label{sec:oracleresults}

To determine whether unsuccessful constituent recovery resulted from the
physical disappearance of a fault signature, a within-pair oracle analysis
was first performed. For each directed compound relation $A|B$,
measurements from the compound state $A{+}B$ were compared with the
corresponding single-$B$ reference while retaining the same constituent of
interest $A$.

All 18 directed constituent--partner tests exhibited strong within-pair
discriminability. The condition-held-out within-pair AUC ranged from 0.9167
to 1.000, with a median of 1.000. Thus, the target compound and its
corresponding partner-only reference remained strongly separable in the
broader physics-feature space. However, the recording-sensitivity control
showed that the two available recordings of the same bearing-outer fault
under the same operating condition were themselves nearly perfectly
separable ($\mathrm{AUC} = 0.998$). High within-pair AUC therefore cannot
be attributed uniquely to the additional constituent and may partly reflect
recording-specific nuisance variation.

A complementary same-date control reduced recording-date mismatch. Eight
pairs of compound states acquired on the same date shared one constituent
and all 12 operating conditions, yielding 16 directed matched comparisons.
Leave-one-condition-out AUCs ranged from 0.972 to 1.000. In particular,
static eccentricity in the bearing-outer context remained perfectly
separable when the alternative same-date compound contained either broken
bar or winding with the same bearing-outer partner ($\mathrm{AUC} = 1.000$
in both comparisons). Because these reference compounds contain a different
non-target constituent, the controls do not causally isolate the target
mechanism; nevertheless, they show that the observed compound-state
discriminability cannot be explained solely by recording-date mismatch.

Accordingly, within-pair detectability does not establish that the same
evidence is transferable across compound partners. Cross-partner
experiments were therefore performed by learning constituent $A$ from other
compound contrasts, such as $A{+}C$ versus $C$, and evaluating the
resulting detector on $A{+}B$ versus $B$.

Under seen operating conditions, the primary cross-partner protocol
achieved a median AUC of 0.905 across the 18 directed tests, with a mean
AUC of 0.815. Thirteen of the 18 directed tests achieved an AUC of at least
0.80, and 15 achieved an AUC of at least 0.70. When the target operating
condition was also withheld, the median AUC remained high at 0.884 and the
mean AUC was 0.803.

These aggregate values indicate that constituent evidence is often
transferable across partners. Nevertheless, the transfer behavior was
highly heterogeneous. Selected examples are reported in
Table~\ref{tab:transfer}. The complete pattern of directional cross-partner
transfer is shown in Figure~\ref{fig:heatmap}.

\begin{table*}[!t]
\caption{Selected directed cross-partner transfer results. The within-pair
AUC measures detectability of constituent $A$ in the target pair, whereas
the cross-partner AUC measures transfer of evidence learned from other
compound partners.}
\label{tab:transfer}
\centering
\begin{tabular}{l l l c c c}
\toprule
Target composition & Constituent $A$ & Partner $B$ &
\makecell{Within-pair\\AUC} & \makecell{Cross-partner\\AUC (seen)} &
\makecell{Cross-partner\\AUC (unseen)} \\
\midrule
Bearing inner $+$ bearing outer        & Bearing outer        & Bearing inner        & 1.000 & 1.000 & 1.000 \\
Bearing inner $+$ dynamic eccentricity & Bearing inner        & Dynamic eccentricity & 1.000 & 0.944 & 0.931 \\
Bearing inner $+$ broken bar           & Broken bar           & Bearing inner        & 1.000 & 0.893 & 0.851 \\
Bearing inner $+$ dynamic eccentricity & Dynamic eccentricity & Bearing inner        & 1.000 & 0.711 & 0.653 \\
Bearing outer $+$ static eccentricity  & Static eccentricity  & Bearing outer        & 0.917 & 0.000 & 0.000 \\
Bearing outer $+$ winding              & Winding              & Bearing outer        & 0.917 & 0.333 & 0.340 \\
\bottomrule
\end{tabular}
\end{table*}

\begin{figure*}[!t]
\centering
\includegraphics[width=\textwidth]{fig3_transfer_heatmap}
\caption{Directional cross-partner transfer AUC for the primary
other-partners transfer protocol. Rows denote the constituent $A$ to be
recovered and columns denote partner $B$ in the target compound. Values
below 0.5 indicate below-chance fixed-orientation ranking and do not by
themselves establish statistically supported polarity reversal. $n_C$
denotes the number of eligible compound training contexts available for the
recovered constituent.}
\label{fig:heatmap}
\end{figure*}

The contrast between within-pair and cross-partner performance was
particularly pronounced for static eccentricity in the presence of the
bearing-outer fault. Within the target pair, static eccentricity remained
strongly detectable, with $\mathrm{AUC}^{\mathrm{within}} = 0.9167$.
However, evidence learned from other eligible compound contexts produced a
fixed-orientation directional AUC of
\begin{equation}
\mathrm{AUC}_{\mathrm{dir}} = 0.000
\end{equation}
under both seen and unseen operating conditions. The corresponding
orientation-free separability was
\begin{equation}
\mathrm{AUC}_{\mathrm{sep}} = 1.000,
\end{equation}
indicating perfect rank separability in the opposite orientation. This was
the only directed constituent--partner relation classified as a
statistically supported score-orientation reversal after multiplicity
correction, with Holm-adjusted $p = 0.0136$ under seen conditions and
$p = 0.0108$ under unseen conditions. The statistical test establishes a
reproducible reversal of the fixed-score ordering, but it does not
establish that the reversal is caused exclusively by a physical inversion
of the static-eccentricity signature. Given the strong recording-identity
separability observed in the nuisance control, partner-dependent
constituent reorientation and acquisition-related variation remain
plausible, non-exclusive contributors to this result.

Winding in the bearing-outer composition also produced below-chance
directional AUCs of 0.3333 and 0.3403 under seen and unseen operating
conditions, respectively. However, its orientation-free separability was
only approximately 0.67 and the Holm-adjusted polarity tests were not
significant. This relation was therefore classified as
orientation-ambiguous rather than as a statistically supported
score-orientation reversal.

Taken together, the within-pair and recording-sensitivity controls show
that compound states can remain strongly discriminable while the
orientation of a score learned from other compound partners may fail to
transfer. Because recording-specific structure is also detectable in the
feature space, this phenomenon is interpreted as context-dependent score
transfer rather than as causal proof of physical constituent-signature
inversion. Figure~\ref{fig:withinvstransfer} directly contrasts within-pair
detectability with directional cross-partner transfer across the evaluated
constituent--partner relations.

\begin{figure}[!t]
\centering
\includegraphics[width=\columnwidth]{fig4_withinpair_vs_transfer}
\caption{Within-pair oracle detectability versus directional cross-partner
transfer. The dashed horizontal line denotes chance-level directional
ranking. The shaded lower region represents below-chance directional
ranking; statistically supported polarity reversal requires the separate
polarity analysis described in Section~\ref{sec:stats}.}
\label{fig:withinvstransfer}
\end{figure}

\subsection{Physical-Signature Heterogeneity and Directional Failure
Modes}
\label{sec:heterogeneity}

Cross-partner transfer was strongly directional. For example, in the
bearing-outer plus static-eccentricity compound, evidence for bearing outer
transferred with an AUC of 1.000, whereas evidence for static eccentricity
in the opposite direction yielded a directional AUC of 0.000.
Transferability was therefore not a symmetric property of the compound
pair.

The feature-level preservation analysis further showed that multivariate
score-orientation reversal should not be equated with universal inversion
of the underlying physical descriptors. For static eccentricity in the
bearing-outer context, the median preservation ratio was 0.760, while
approximately 18.5\% of the eligible curated feature-condition observations
exhibited feature-level sign reversal and 31.4\% exhibited severe
suppression. Thus, most individual physical descriptors did not reverse
orientation even though the resulting multivariate cross-partner ranking
was completely inverted.

Conversely, strong feature suppression did not necessarily imply
score-orientation reversal. Bearing-outer evidence in the bearing-inner
compound showed a median preservation ratio of approximately 0.305 and a
severe-suppression rate of approximately 58\%, yet its cross-partner
directional AUC remained 1.000. Feature-level suppression and multivariate
score-orientation reversal therefore represent distinct forms of
partner-dependent behavior. Figure~\ref{fig:preservation} illustrates this
feature-level heterogeneity for static eccentricity under the two
bearing-partner contexts.

\begin{figure}[!t]
\centering
\includegraphics[width=\columnwidth]{fig5_preservation}
\caption{Feature-level preservation of selected static-eccentricity
signatures under bearing-inner and bearing-outer partner contexts. Negative
values indicate sign reversal, values near one preservation, and values
above one amplification.}
\label{fig:preservation}
\end{figure}

For the locked logistic-regression probe, the false-addition values in
Table~\ref{tab:crossed} should not be interpreted as conventional
false-alarm probabilities. The metric represents the mean number of
inactive constituent labels incorrectly activated per compound prediction.
Nevertheless, values of approximately 0.64--0.73 indicate a non-negligible
tendency toward over-diagnosis. In a practical condition monitoring system,
these additional constituent activations could increase inspection burden
even when one or both true mechanisms are successfully identified.
Deployment-oriented extensions should therefore consider cost-sensitive
threshold calibration or asymmetric decision penalties that explicitly
trade constituent recall against false additions. The classifier-family
robustness analysis further showed that this over-diagnosis tendency is not
fixed by the constituent representation itself: random forest and the
neural baseline reduced mean FA/run to approximately 0.26 while retaining
mean constituent recall near 0.66.

The observed cross-partner score reorientation also suggests a potential
role for representation or normalization strategies that explicitly promote
constituent consistency across compound contexts. Possible directions
include robust condition-aware normalization learned exclusively from
training data, partner-insensitive spectral reference schemes, or
representation constraints that align evidence for the same constituent
across different compound partners. However, the present preservation
analysis shows that the observed failure cannot be reduced to a uniform
inversion of the underlying physical features. Such adaptation strategies
should therefore be regarded as hypotheses for future investigation rather
than as established remedies for score-orientation reversal.

A further limitation is the recording structure of the public dataset.
Acquisition timestamps recovered from the released file names showed that
distinct recordings of the same fault state can be strongly separable in
the physics-feature space. Although same-date shared-partner controls
retained near-perfect discriminability, recording date is only a coarse
proxy for experimental session, and the available release does not provide
replicated same-state measurements sufficient to separate fault, assembly,
sensor-mounting, thermal, and recording-session effects. The cross-partner
results should therefore be interpreted as evidence of context-dependent
score transfer rather than as causal proof of physical signature inversion.
Absolute constituent-set reconstruction was additionally sensitive to the
classifier family: nonlinear tree and neural baselines substantially
reduced false additions and improved exact-match recovery, while the
descriptive advantage of recombination over complete compound-context
deprivation persisted across all four classifier families.

Taken together, the experiments identify two distinct limitations of
constituent-level compound-fault generalization. First, reduced
compound-context exposure became measurably detrimental when evaluation
also required generalization to an unseen operating condition. Second,
cross-partner transfer could be strongly directional even when within-pair
constituent evidence remained highly detectable. The latter included one
statistically supported case of complete score-orientation reversal, while
the available recording structure prevents attributing that reversal
uniquely to a physical signature inversion.
